\section{Introduction}
\label{sec:introduction}

\subsection{The Smallest Useful Splitter}

BISECT has several ways to split a bisection-tree node: METIS, GeoSection,
CVD, simulated annealing, flow, exact methods, and others. \bfsg{} is the
opposite end of that spectrum. It asks what happens if a split is built with
only the tract graph, tract populations, and a deterministic seed stream.

At each node, \bfsg{} does four things:
\begin{enumerate}
\item choose one seed by population-weighted sampling;
\item choose the second seed at maximum BFS distance from the first;
\item grow two frontier regions, always expanding the side with smaller
  current population first; and
\item apply a local boundary rebalance pass.
\end{enumerate}

This makes \bfsg{} a useful baseline because it is easy to explain and easy to
audit. It does not claim to optimize edge cut, Polsby-Popper, Reock, partisan
fairness, or statutory compliance. It is a deliberately small contiguous
growth rule that can be compared against stronger algorithms.

\subsection{Show the Growth}

The core idea should be visible, not merely named. In the schematic below, the
two seeds start on opposite graph ends. The smaller side expands first. The
frontier moves only through adjacency.

\begin{center}
\begin{tabular}{c@{\qquad}c@{\qquad}c}
\textbf{1. Seeds} & \textbf{2. Frontier} & \textbf{3. Split} \\
\begin{tabular}{|c|c|c|c|c|}
\hline
S0 & . & . & . & S1 \\ \hline
. & . & . & . & . \\ \hline
. & . & . & . & . \\ \hline
\end{tabular}
&
\begin{tabular}{|c|c|c|c|c|}
\hline
0 & 0 & . & 1 & 1 \\ \hline
0 & . & . & . & 1 \\ \hline
. & . & . & . & . \\ \hline
\end{tabular}
&
\begin{tabular}{|c|c|c|c|c|}
\hline
0 & 0 & 0 & 1 & 1 \\ \hline
0 & 0 & 1 & 1 & 1 \\ \hline
0 & 0 & 1 & 1 & 1 \\ \hline
\end{tabular}
\end{tabular}
\end{center}

This is why \bfsg{} belongs in the atlas and paper set: it gives a reader a
plain visual model for what ``region growing'' means before the more complex
optimization families enter the story.

\subsection{Implemented Scope}

The current Rust implementation is a recursive bisection structure:
\texttt{run\_all\_splits\_bfs} walks the standard \texttt{BisectionTree} and
calls \texttt{split\_subgraph\_bfs} at each internal node. It does not grow
$k$ districts simultaneously. Odd-seat nodes inherit the same broad issue
seen in CVD: the bisection tree records child seat counts, but this splitter's
local rebalance is still oriented around a half-population split.

The implementation uses two seed functions:
\begin{itemize}
\item \texttt{BFS\_NODE\_} derives the per-node seed from base seed and node path.
\item \texttt{BFS\_SEED\_} is then used inside the split to seed the
  population-weighted first-seed sample.
\end{itemize}

That two-stage seed derivation should be preserved in any audit package:
same base seed plus same node path implies same seed placement and same split.

\subsection{What This Paper Claims}

This paper makes three modest claims.

\paragraph{Algorithmic claim.}
\bfsg{} is a deterministic, seed-auditable bisection primitive once the base
seed and node path are fixed.

\paragraph{Reader claim.}
\bfsg{} is the easiest construction algorithm to teach: seed, grow, rebalance.
It should be used as the first visual baseline in explanations of the
structure-layer family.

\paragraph{Evidence claim.}
The code has useful L0/L1 tests for validity, determinism, contiguity, and CLI
parsing. Statewide quality and MCMC warm-start claims remain benchmark targets
until archived run packages exist.
